\SourceLeaf{611}\begin{center}[4.]\end{center}
After a hundred and fifty years, we still know next to nothing about \textit{Geometria Situs} and have scarcely anything to show for it. \textsc{Leibnitz} had an inkling of it, and just two geometers, \textsc{Euler} and \textsc{Vandermonde}, were afforded a faint glimpse into it.

One of the chief problems in the \textit{borderland} between \textit{Geometria Situs} and \textit{Geometria Magnitudinis} will be to count how often two closed or infinite lines wind around each other.

Take $x,y,z$ as the coordinates of an unspecified point on the first line and $x',y',z'$ as those on the second, and set
\[
\int\int\frac{\begin{gathered}
(x'-x)(\mathrm dy\,\mathrm dz'-\mathrm dz\,\mathrm dy')\\
+(y'-y)(\mathrm dz\,\mathrm dx'-\mathrm dx\,\mathrm dz')\\
+(z-z')(\mathrm dx\,\mathrm dy'-\mathrm dy\,\mathrm dx')
\end{gathered}}{[(x'-x)^2+(y'-y)^2+(z'-z)^2]^{\frac{3}{2}}}=V
\]
Then, integrating along both lines, the value of the integral is
\[
=4m\pi
\]
where $m$ counts the turns of the two lines around each other.

The value is symmetric. Interchanging the two lines leaves it unchanged.\quad 1833.\ Jan.\ 22.

\begin{center}[5.]\par\medskip\textit{Law of the Galvanic Current:}\end{center}
1.\quad The positive current is the one that flows, at the point of contact with the water, in the direction zinc, water, copper.

2.\quad Let $RR'$ be a current element in which the positive current runs from $R$ to $R'$, and let $P$ be a point containing an element of positive northern magnetic fluid.

The current element $RR'=\mu$ then exerts a force on $P$ of magnitude
\[
=\frac{\mu\,.\,\sin R'RP}{(RR')^2}
\]
Its direction $PQ$ is perpendicular to the plane through $P,R,R'$ and points downward if $P$ lies to the right of $RR'$ or $R'$ lies to the right of $PR$.
\par
\SourceLeaf{612}
3.\quad The action of a closed current $RR'R''R$ can be matched by magnetic action in the following way. Let $RR'R''R$ bound an arbitrary surface, and imagine northern magnetism spread over it in any distribution, with density $=\delta$. At each point of the surface erect an infinitesimal normal, directly proportional to the intensity of the current and inversely proportional to $\delta$. Let it point upward or downward according as the current, in going around the surface, has it on its right or its left. The endpoints of these normals lie on a second surface. Spread over it and over its parts exactly as much southern magnetism as there is on the other surface and its corresponding parts. These two magnetic surfaces are equivalent to that galvanic current at every point outside them.

